\documentclass[../../main.tex]{subfiles}
\begin{document}

{\bf [解]}

与えられた定積分
\begin{align}
 S(a,t) &= \int_0^a\left|e^{-x}-\dfrac{1}{t}\right|dx \label{eq:1}
\end{align}
について考える．被積分関数は$y=e^{-x}$と直線$y=\dfrac{1}{t}$の差として図形的に理解できる．
$y=e^{-x}$は$[0,a]$で$1$から$e^{-a}$まで単調減少することに注意すると，グラフは以下のようになる．


\paragraph{(1)}
$0 \le x \le a$において$e^{-a} \le e^{-x} \le 1$だから，$\dfrac{1}{t}$とこの大小関係によって以下のように場合分けする．

\subparagraph{$1/t\ge1$（$t\le1$）の時} 
$[0,a]$上で常に$e^{-x} \le \dfrac{1}{t}$だから
\begin{align}
S(a,t)=\int_0^a\left(\dfrac{1}{t}-e^{-x}\right)dx=\dfrac{a}{t}+e^{-a}-1
\end{align}
である．
これは$(1/t)$の1次関数で係数$a>0$だから，$(1/t)$の単調増加関数すなわち$t$の単調減少関数である．

\subparagraph{$1/t\le e^{-a}$（$t\ge e^a$）の時}
$[0,a]$上で常に$e^{-x}\ge \dfrac{1}{t}$だから
\begin{align}
S(a,t)=\int_0^a\left(e^{-x}-\dfrac{1}{t}\right)dx=(1-e^{-a})-\dfrac{a}{t}
\end{align}
である．
これは$(1/t)$の単調減少関数すなわち$t$の単調増加関数である．
\casesend

さて$S(a,t)$は明らかに連続だから，以上の場合分けを踏まえて$e^{-a}\le1/t\le1$（$1\le t\le e^a$）の範囲で最小値をとる．
この範囲では，
\begin{align}
 e^{-\alpha}=\dfrac{1}{t} \label{eq:2}
\end{align}
をみたす$\alpha\in[0,a]$がただ1つ存在して$x<\alpha$で$e^{-x}>1/t$，$x>\alpha$で$e^{-x}<1/t$だから$S(a,t)$は
\begin{align}
\begin{split}
S(a,t)
 &=\int_0^\alpha\left(e^{-x}-\dfrac{1}{t}\right)dx+\int_\alpha^a\left(\dfrac{1}{t}-e^{-x}\right)dx \\
 &=\left[-e^{-x}-\dfrac{x}{t}\right]_0^\alpha+\left[\dfrac{x}{t}+e^{-x}\right]_\alpha^a \\
 &=\left(1-e^{-\alpha}-\dfrac{\alpha}{t}\right)+\left(\dfrac{a}{t}+e^{-a}-\dfrac{\alpha}{t}-e^{-\alpha}\right) \\
 &=1+e^{-a}-2e^{-\alpha}+\dfrac{a-2\alpha}{t} \quad\cdots\text{①}
\end{split} \label{eq:3}
\end{align}
と書ける．\cref{eq:2}より$\alpha = \log t$に注意して\cref{eq:3}から$\alpha$を消去すると
\begin{align}
 S(a,t)
  &= 1+e^{-a}+\dfrac{-2+a-2\log t}{t} \label{eq:4}
\end{align}
である．以下$a$を固定した時の$S(a,t)$の最小値を求めよう．
$t$で微分して
\begin{align}
\dfrac{dS}{dt}
  &= \dfrac{\left(-\dfrac{2}{t}\cdot t -\left(-2+a-2\log t\right) \right)}{t^2} \\
  &= \dfrac{- a + 2\log t}{t^2} 
\end{align}
である．従って増減表は以下のようになる．

\begin{table}[htb]
  \centering
\caption{$S$の増減表}
\begin{tabular}{c|ccccc}
$t$ & $1$ & & $e^{a/2}$ & & $e^a$ \\
\hline
$S'$ & & $-$ & $0$ & $+$ & \\
\hline
$S$ & & $\searrow$ & & $\nearrow$ &
\end{tabular}
\end{table}

従って$S$は$t=e^{a/2}$で最小値$m(a)$をとるから\cref{eq:4}に代入して
\begin{align}
  m(a)=\min S\left(e^{a/2}\right) = 1+e^{-a}-2e^{-a/2}=\left(1-e^{-a/2}\right)^2
\end{align}
である．

\paragraph{(2)}
簡単のため
\begin{align}
 p = \dfrac{-a}{2}
\end{align}
とおくと$a\to0$で$p\to0$であって
\begin{align}
\dfrac{m(a)}{a^2}=\dfrac{\left(1-e^{-a/2}\right)^2}{a^2}=\dfrac{\left(e^{p}-1\right)^2}{(-2p)^2}=\dfrac{1}{4}\left(\dfrac{e^p-1}{p}\right)^2
\end{align}
となる．ここで$e^{x}$の微分を考えると
$\dfrac{e^p-1}{p}\to1$（$p\to0$）だから求める極限値は
\begin{align}
\lim_{a\to0}\dfrac{m(a)}{a^2}=\dfrac{1}{4}
\end{align}
である．

\newpage

\end{document}
